\documentclass[10pt, twocolumn, letterpaper]{article}

\usepackage[utf8]{inputenc}
\usepackage[T1]{fontenc}
\usepackage{mathptmx} % Times New Roman clone for standard journal feel
\usepackage{geometry}
\geometry{left=1in, right=1in, top=1in, bottom=1in, columnsep=0.25in}
\usepackage{titlesec}
\usepackage{cite}
\usepackage{hyperref}

\titleformat{\section}{\normalsize\bfseries\uppercase}{\thesection.}{1em}{}
\titleformat{\subsection}{\normalsize\itshape}{\thesubsection.}{1em}{}

\title{\textbf{The Forcing Function of Autonomy:\\Structural Accountability and the Right to Exit}}
\author{Justin Bogner\thanks{Independent Scholar. Contact: justin@example.com}}
\date{\vspace{-5ex}}

\begin{document}

\twocolumn[
  \begin{@twocolumnfalse}
    \maketitle
    \begin{abstract}
    The expansion of Medical Aid in Dying (MAiD) to non-terminal populations frequently encounters the disability rights critique, which contends that systemic deprivation---such as poverty and lack of social support---coerces marginalized individuals into choosing death. This paper acknowledges the validity of the critique but rejects the conventional protectionist remedy of categorically denying MAiD eligibility. Prohibiting access to a safe exit does not resolve the underlying deprivation; it merely traps the marginalized in a state of indefinite suffering to shield the state from political embarrassment. Instead, this paper proposes the ``Sovereignty Model'' coupled with a ``Structural Accountability Mechanism.'' By legally decoupling MAiD eligibility from state-adjudicated suffering and requiring mandatory administrative reporting of the socioeconomic drivers behind MAiD requests, the framework transforms an individual right to exit into a political forcing function. This mechanism exposes the state's failure to provide adequate care, leveraging the profound moral weight of voluntary death to catalyze social and economic investment without sacrificing the individual's bodily autonomy.
    \vspace{2em}
    \end{abstract}
  \end{@twocolumnfalse}
]

\section{Introduction}

The debate surrounding the expansion of Medical Aid in Dying (MAiD) for non-terminal conditions has exposed a profound tension between individual autonomy and systemic inequality. Disability rights advocates---including the UN Special Rapporteur on the Rights of Persons with Disabilities---argue that offering MAiD in an environment characterized by severe socioeconomic deprivation amounts to a form of passive eugenics \cite{quinn2021}. When adequate housing, accessible care, and sufficient disability income are unavailable, a request for MAiD cannot be construed as a free choice; it is coerced by systemic abandonment.

In response, the dominant legislative strategy has been a retreat into a ``Protectionist Model,'' which attempts to insulate the vulnerable by enforcing stringent clinical gatekeeping and denying eligibility to populations perceived to be uniquely susceptible to systemic coercion. 

This paper argues that the protectionist remedy is legally incoherent and morally unacceptable. Trapping a marginalized person in a state of suffering they deem intolerable does not solve the underlying deprivation. Instead, it creates a mechanism that this paper terms the ``Sovereignty Model,'' featuring a Structural Accountability Mechanism designed to use individual autonomy as a political forcing function.

\section{The Fallacy of the Protectionist Remedy}

The disability rights critique correctly diagnoses the disease (systemic deprivation) but prescribes a fatal remedy (categorical prohibition of exit).

Consider the logic of the protectionist state: \textit{We recognize that we have failed to provide you with a livable income, adequate housing, or necessary care. Because our failures have rendered your life intolerable, we must conclude that your request to exit this life is not an expression of authentic autonomy, but a symptom of our neglect. Therefore, to protect you from the consequences of our neglect, we forbid you from exiting, forcing you to continue suffering under the conditions we refuse to fix.}

This is not protection; it is institutional entrapment \cite{stefan2016}. To solve state-created danger by criminalizing the escape route is to compound systemic failure with paternalistic cruelty. The state leverages its own negligence to disenfranchise the individual from their most fundamental liberty interest: the right to bodily autonomy and informed refusal \cite{beauchamp2019}.

\section{The Structural Accountability Mechanism}

If prohibition is not the answer, how can jurisprudence and public policy reconcile the right to exit with the reality of socioeconomic coercion? The answer lies in the ``Structural Accountability Mechanism'' (SAM).

The SAM decouples MAiD eligibility from suffering and re-grounds it entirely in decisional capacity. It acknowledges that a capacitated individual's reason for requesting MAiD may be fundamentally structural (e.g., an inability to afford out-of-pocket home care). However, rather than disqualifying the applicant, the SAM mandates a rigorous, legally formalized process of socioeconomic driver identification and reporting.

\subsection{Mandatory Driver Reporting}

Under the SAM, prior to any MAiD provision, an independent, unaffiliated social worker must conduct a coercion and resource screening. If the applicant cites lack of housing, insufficient income, or denied insurance coverage as primary or contributing drivers, the state is legally obligated to immediately offer those specific resources.

If the resources are rejected by the applicant, the state's obligation is discharged. However, if the resources are simply \textit{unavailable} (due to geographic lack of services or financial inaccessibility), the state must document this exact failure on a mandatory reporting form.

\section{The Forcing Function of Autonomy}

The true power of the SAM is not merely in documenting failure, but in political visibility. The conventional protectionist model hides state failure by burying it in private clinical encounters, allowing the state to quietly deny both the necessary support and the right to die.

By contrast, the SAM creates a public administrative record. When a state's Department of Health is required to publish an annual report revealing that $X$ number of citizens chose to end their lives specifically because the state failed to provide affordable housing or adequate disability support, the political equilibrium fundamentally shifts.

The individual's exercise of autonomy becomes a ``forcing function.'' It weaponizes the profound moral and public shock of voluntary death to expose the hypocrisy of the state. It strips away the comforting illusion of protectionism, forcing legislatures to directly confront the lethal consequences of their austerity measures. Over time, the political cost of these documented deaths generates the exact legislative pressure required to fund the social determinants of health. 

\section{Conclusion}

The disability rights movement is correct that systemic deprivation coerces the marginalized, but the answer cannot be the paternalistic restriction of fundamental rights. 

The Sovereignty Model offers a radical alternative: it honors the individual's ultimate jurisdiction over their own life, while implementing a Structural Accountability Mechanism that forces the state to legally document and politically own its failures. By aligning bodily autonomy with systemic visibility, we can create a legal framework that is both deeply libertarian in its respect for the individual and powerfully progressive in its demand for social justice.

\begin{thebibliography}{9}
\small
\bibitem{beauchamp2019} Beauchamp, T.L., \& Childress, J.F. (2019). \textit{Principles of biomedical ethics} (8th ed.). Oxford University Press.
\bibitem{quinn2021} Quinn, G. (2021). \textit{Report of the Special Rapporteur on the rights of persons with disabilities}. United Nations General Assembly.
\bibitem{stefan2016} Stefan, S. (2016). \textit{Rational suicide, irrational laws: Examining current approaches to suicide in policy and law}. Oxford University Press.
\end{thebibliography}

\end{document}
